\chapter{Anomaly Detection}
\label{ch:anomaly}

\chapterguide%
  {\chapref{ch:unsupervised}, distributions and quantiles from \chapref{ch:foundations}, plus clustering and density ideas from \chapref{ch:clustering}.}
  {score anomalies, compare detectors, set review thresholds, and evaluate alerts.}

\term{Anomaly detection} ranks departures from a reference pattern, such as fraud, sensor failures, or entry errors.

\section{Anomaly detection}

Two closely related setups share the name:

\begin{mltable}
\begin{tabularx}{\linewidth}{@{}>{\bfseries\leavevmode\color{MLNavy}}p{0.22\linewidth}p{0.36\linewidth}X@{}}
\toprule
\rowcolor{MLPaleBlue!92}
Setup & Training data & Goal \\
\midrule
Outlier detection & The training data itself is contaminated with anomalies. & Find the anomalous observations already present in the data. \\
Novelty detection & The training data contains clean ``normal'' behavior. & Flag future observations that deviate from the learned normal pattern. \\
\bottomrule
\end{tabularx}
\end{mltable}

Because anomalies are rare and costly to label, many methods learn a reference pattern and score departures from it.

\subsection{Statistical rules: z-scores and IQR fences}

For a single roughly bell-shaped feature, the \term{z-score} measures how many standard deviations an observation sits from the mean:
\[
  z=\frac{x-\bar{x}}{s},
\]
with $|z|>3$ a customary alarm threshold --- under a normal distribution, only about $0.3\%$ of values land that far out.

\begin{workedexample}
A temperature sensor averages $\bar{x}=50$ with standard deviation $s=5$. A reading of $68$ has $z=(68-50)/5=3.6$: flagged. A reading of $57$ has $z=1.4$: unremarkable.
\end{workedexample}

Outliers can inflate the mean and standard deviation, masking themselves. Robust alternatives use the median and scaled median absolute deviation. \term{IQR fences} use quartiles instead (\chapref{ch:statistics}):
\[
  [\,Q_1-1.5\operatorname{IQR},\ Q_3+1.5\operatorname{IQR}\,],
\]
For $Q_1=20$ and $Q_3=30$ minutes, IQR fences are 5 and 45; a 60-minute delivery is flagged without assuming normality.

Per-feature checks can miss unusual combinations of individually plausible values. Multivariate methods capture these relationships.

\subsection{Distance and local density}

Distance to the $k$-th nearest neighbor measures isolation, but naturally sparse groups can look anomalous. \term{Local outlier factor (LOF)} compares density with neighboring densities: near~1 is locally typical; much larger values indicate local isolation.

\subsection{Model-based scores}

Fit a Gaussian or mixture to reference data and flag low-likelihood cases (\chapref{ch:clustering}). A \term{one-class SVM} instead learns a boundary around the reference population.

\subsection{Isolation forests}
\label{sec:isolation-forest}

Isolation forests build trees with random feature thresholds and measure how quickly each observation becomes isolated.

\begin{intuition}
Crowded points need many random splits to isolate; isolated points need fewer. Short normalized paths therefore indicate greater anomaly.
\end{intuition}

Isolation forests need no distance metric. \term{Contamination} sets a threshold, not the true anomaly rate; Lab~5 ranks raw scores instead.

\subsection{Evaluating and using anomaly detectors}

\begin{itemize}
  \item Evaluate reviewed top-$k$ precision and, with labels, precision--recall. Lab~5's top 100 alerts have 20\% precision against 1.5\% prevalence; AUC 0.95 measures a different aspect.
  \item For distance- and kernel-based detectors, apply the scaling rule from \secref{sec:scaling}.
  \item Choose thresholds from consequences --- how many alarms per day can be investigated? --- not from statistical custom.
  \item Anomalies may be errors or important rare events. Investigate before deleting.
\end{itemize}

\section{The shape of an anomaly problem}

Anomaly detection differs from classification in one important way: you usually cannot see many examples of the thing you are looking for.

\begin{mlfigure}
\begin{tikzpicture}[font=\scriptsize]
\begin{scope}
\draw[draw=MLMuted,fill=white] (0,0) rectangle (4.2,2.9);
\foreach \p in {(1.5,1.4),(1.8,1.6),(1.6,1.1),(2.0,1.3),(1.3,1.7),(2.2,1.6),(1.9,0.9),
                (1.4,1.25),(2.1,1.1),(1.7,1.45),(2.3,1.35),(1.55,0.85),(2.05,1.65)}
  \fill[MLBlue] \p circle (2.4pt);
\draw[MLTeal,thick,dashed] (1.8,1.32) circle (0.95);
\fill[MLRed] (3.5,2.4) circle (2.8pt);
\fill[MLRed] (0.5,0.4) circle (2.8pt);
\node[color=MLRed,anchor=west] at (3.0,2.72) {anomalies};
\node[color=MLTeal] at (1.8,-0.35) {the "normal" region};
\node[font=\small\bfseries\color{MLNavy}] at (2.1,3.25) {Distance / density view};
\end{scope}
\begin{scope}[xshift=5.6cm]
\draw[draw=MLMuted,fill=white] (0,0) rectangle (4.6,2.9);
\draw[-{Stealth[length=2mm]},MLMuted] (0.3,0.45) -- (4.35,0.45);
\node[color=MLMuted,anchor=north] at (2.3,0.35) {anomaly score};
\draw[MLBlue,very thick] plot[smooth,domain=0.4:4.3,samples=60]
  (\x, {0.45+2.0*exp(-((\x-1.3)^2)/0.28)});
\draw[MLRed,very thick,dashed] (3.05,0.45) -- (3.05,2.6);
\node[color=MLRed,anchor=west,align=left] at (3.15,2.2) {threshold\\you choose};
\node[color=MLGreen,anchor=east] at (2.2,1.05) {normal};
\node[color=MLRed,anchor=west] at (3.3,0.85) {flagged};
\node[font=\small\bfseries\color{MLNavy}] at (2.3,3.25) {Score-and-threshold view};
\end{scope}
\end{tikzpicture}
\end{mlfigure}

\begin{keyidea}
Choose thresholds or review batches for actual capacity and error costs; a high score is not a verdict.
\end{keyidea}

\begin{analogy}
Anomaly detection learns a reference pattern and flags departures, without requiring examples of every failure.
\end{analogy}

\section{From anomaly score to investigation}

Define the review process alongside the scoring model.

\begin{mltable}
\begin{tabularx}{\linewidth}{@{}>{\bfseries\leavevmode\color{MLNavy}}p{0.35\linewidth}X@{}}
\toprule
\rowcolor{MLPaleBlue!92}
Question & Practical consequence \\
\midrule
How many cases can be reviewed? & Choose a top-$k$ or score threshold that respects review capacity. \\
What happens after a flag? & Define whether the case is blocked, queued, inspected, or merely logged. \\
Which false alert is costly? & Track precision among reviewed cases and the burden on users or analysts. \\
Can normal behavior change? & Monitor score distributions and retrain or recalibrate when the reference population drifts. \\
Can labels arrive later? & Use confirmed outcomes to evaluate ranking quality and eventually consider a supervised model. \\
\bottomrule
\end{tabularx}
\end{mltable}

\begin{modelcard}{Anomaly Detection - Quick Guide}
\begin{tabularx}{\linewidth}{@{}>{\bfseries\leavevmode\color{MLNavy}}p{0.25\linewidth}X@{}}
Anomaly detection asks & Which observations are unusually inconsistent with a reference pattern? \\
Useful methods & Robust statistical rules, distance/local-density methods, isolation forest, and one-class models. \\
Output & An anomaly score or ranking; a later policy converts that score into an alert. \\
Evaluation & Precision among reviewed cases, PR AUC when labels exist, stability, and downstream review value. \\
Main risks & Rare does not always mean harmful; distribution shift can make yesterday's normality model obsolete. \\
\end{tabularx}
\end{modelcard}

\begin{quickcheck}
\begin{enumerate}
  \item Why is an anomaly score not automatically a probability of fraud, failure, or error?
  \item How does a fixed review capacity turn anomaly detection into a ranking problem?
  \item What is the difference between detecting outliers in the training sample and novelty detection for future observations?
\end{enumerate}
\end{quickcheck}
